\section{Related Work}
\label{sec:related}
\paragraph{Pill identification and low-shot retrieval.}
Consumer-to-reference matching was established as a medication recognition task by the NLM Pill Image Recognition Challenge~\citep{yaniv2016nlm}, followed by fine-grained pill identification using learned convolutional features and comparisons with handcrafted descriptors~\citep{wong2017pill}. These studies demonstrate the usefulness of appearance while exposing the difficulty of separating similar products under image variation within a product. This setting is developed further by ePillID, whose strongest low-shot baseline combines bilinear features with multi-head metric learning~\citep{usuyama2020epillid}. Its analysis of confusing classes is relevant to personal galleries containing medicines with similar appearances, although ranked retrieval from product references leaves gallery membership rejection as a separate decision. PharmaScan evaluates both tasks and distinguishes results under strict and official protocols because access to references for test identities changes the information available during representation learning.

\paragraph{Multiple visual cues and feature fusion.}
Complementary visual cues are useful when discriminating pill characteristics are not equally visible across images. Ling et al. introduce CURE images with varied acquisition conditions, a segmentation network and a multi-stream recognition model trained using Batch All followed by Batch Hard mining~\citep{ling2020fewshot}. Their combination of pill cues and difficult training examples relates to the broader problem of recognition beyond familiar photographs. In PharmaScan, convolutional neural network (CNN) features~\citep{simonyan2015vgg} are combined with vision transformer (ViT) features~\citep{dosovitskiy2021vit}, bringing together local spatial structure and self-attention between image patches~\citep{vaswani2017attention}. These complementary spatial assumptions also connect the architecture with the course treatment of vision transformers~\citep[p.~36]{md6329lecture09}. Channel and spatial attention are provided by the Convolutional Block Attention Module (CBAM)~\citep{woo2018cbam}, while capsule routing aggregates branch features through agreement~\citep{sabour2017capsules}. The inherited architecture supplies the representation in Section~\ref{sec:arch}, whose performance is evaluated within the complete verification workflow.

\paragraph{Metric learning for enrolment.}
Similarity representations separate enrolment from network training by mapping identity information to a space used for retrieval. FaceNet learns this geometry through triplet supervision~\citep{schroff2015facenet}, while CosFace and ArcFace impose margins in cosine and angular space, respectively~\citep{wang2018cosface,deng2019arcface}. These objectives relate identity supervision to distances between examples, connecting with the course discussion of embedding geometry and matching views~\citep[pp.~38, 47--49]{md6329lecture08}. PharmaScan uses ArcFace during training and prototypes during inference, allowing the deployed gallery to contain identities outside the training classifier. Balanced identity sampling~\citep{hermans2017triplet} and evaluation protocol~\citep{musgrave2020reality} are relevant to interpreting such embeddings because optimisation and sampling changes can affect apparent differences between objectives. Accordingly, the sequential analysis in Section~\ref{sec:ablations} evaluates the combined training recipe without isolating an advantage attributable to the loss alone.

\paragraph{Open set decisions and confidence calibration.}
Open set recognition extends discrimination among known identities with a membership decision~\citep{scheirer2013openset}, for which score-based detection provides a baseline~\citep{hendrycks2017baseline}. Membership in PharmaScan is defined by the patient's enrolled gallery, so a medicine familiar to the encoder can still be foreign to the patient. The consequences of false acceptance motivate selection of the operating point on validation data and separate reporting of abstention, consistent with the course treatment of decision thresholds~\citep[p.~92]{md6329lecture02}. Calibration concerns the relationship between reported probabilities and observed correctness: Guo et al. document poorly calibrated network confidence and examine post-processing corrections~\citep{guo2017calibration}. Here, logistic calibration~\citep{platt1999probabilistic} is fitted to catalogue retrieval evidence, while gallery rejection and ambiguity retain their own thresholds. These distinct purposes underlie the decision mechanisms in Sections~\ref{sec:decision} and~\ref{sec:identify}.

\paragraph{Imprint evidence and acquisition transfer.}
Florence-2 represents multiple visual tasks through a common prompted model and supplies the pretrained imprint reader used here~\citep{xiao2024florence2}. Hidden, small or debossed text limits the evidence available to this reader, while adaptation introduces a further distinction between text recovery and verification performance. Low-rank adaptation (LoRA) provides a parameter-efficient mechanism for adapting pretrained models~\citep{hu2022lora}, but an incorrect token can promote a visually similar product even when average recovery improves. The analysis of reader adaptation therefore measures accepted identities alongside recovered text. Acquisition transfer is examined with MEDISEG, which adds instance annotations for Hong Kong products photographed under conditions different from the US reference collection~\citep{mediseg2026,chu2026mediseg}. Changes in scale, scene and product population extend the evaluation beyond retrieval within the source dataset to the transfer of both representation and calibrated decisions.
